\section{Architettura}\label{sec:architettura}

\subsection{Panoramica}
Il sistema è una DApp in cinque parti (tabella~\ref{tab:componenti}). Paziente e medico usano un'applicazione web, che dialoga con il backend per identità e dati clinici e direttamente con la blockchain, tramite il wallet del paziente, per concedere o revocare un consenso.

\begin{table}[H]
\centering\footnotesize
\renewcommand{\arraystretch}{1.2}
\begin{tabular}{@{}p{.17\textwidth}p{.27\textwidth}p{.38\textwidth}p{.1\textwidth}@{}}
\toprule
Componente & Tecnologia & Compito & Stato\\
\midrule
Frontend & React, Vite, WebAuthn del browser & Accesso con passkey, gestione delle proprie passkey e delle sessioni, consultazione delle cartelle e firma dei referti, concessione e revoca dei consensi con il wallet del paziente, registro degli accessi verificabile sulla catena & Completo\\
Backend & Node.js 24, TypeScript, Express 5, ethers 6 & Autenticazione FIDO2, sessioni, controllo degli accessi, cifratura e firma dei referti, scrittura dell'audit, registro per il paziente, cancellazione su richiesta & Completo\\
Database & MongoDB 8.0 & Utenti, credenziali FIDO2, sessioni, profili, referti cifrati, registro cifrato degli accessi & Completo\\
Custodia delle chiavi & AES-256-GCM, ML-KEM-768 + X25519, Argon2id, Shamir & Chiavi dell'archivio cifrate con passphrase, backup a soglia & Completo\\
Smart contract & Solidity 0.8, OpenZeppelin, Hardhat & Registro dei consensi per finalità, registro di audit con accessi d'emergenza, amministrazione con passaggio di consegne & Completo\\
\bottomrule
\end{tabular}
\caption{Componenti del sistema e stato del prototipo.}\label{tab:componenti}
\end{table}

Il backend è l'unico componente che vede i dati clinici in chiaro, e solo per servirli a un utente autorizzato, ma non è l'unica fonte di verità: non può concedere consensi al posto del paziente, perché il contratto accetta solo transazioni firmate dal paziente; non può registrare come autorizzato un accesso senza consenso, perché il registro di audit lo rifiuta; non può riscrivere gli accessi registrati, perché il registro è in sola aggiunta. Compromettere il backend non basta quindi a nascondere un accesso abusivo.

\subsection{Organizzazione secondo il pattern MVC}
Il codice segue il pattern Model-View-Controller, adattato alla presenza di due archivi, uno tradizionale e uno distribuito. Il \textbf{Model} conserva lo stato del sistema dietro tre interfacce:
\begin{itemize}
  \item \texttt{Store}, le operazioni sui dati off-chain, implementata da \texttt{MongoStore} nel server e da \texttt{MemoryStore} nei test;
  \item \texttt{Ledger}, la vista sulla blockchain (leggere un consenso, scrivere una voce di audit), implementata da \texttt{EthersLedger}, che parla con i contratti, e da \texttt{MemoryLedger}, che ne riproduce le regole nei test;
  \item \texttt{Vault}, che custodisce le chiavi dell'archivio e offre solo cifratura e decifratura, senza esporre le chiavi.
\end{itemize}
Sulla catena il Model sono gli smart contract.

La \textbf{View} è il frontend React, che non contiene regole di sicurezza: ogni decisione è presa dal backend o dai contratti. I \textbf{Controller} sono le funzioni del backend che validano la richiesta, applicano le regole (passkey, ruolo, consenso, sessione) e interrogano il Model; le rotte li collegano ai middleware comuni (sessione, limite di frequenza, intestazioni di sicurezza, gestione degli errori).

La separazione serve alla sicurezza: i vincoli importanti, come l'unicità di una credenziale o il consumo atomico di una challenge, stanno in un solo punto e non possono essere dimenticati in una rotta. Inoltre l'applicazione riceve il Model dall'esterno (\texttt{createApp(store, ledger, vault)}), quindi i test eseguono lo stesso codice dei controller che gira nel server.

\subsection{Cosa sta sulla blockchain e cosa no}
Sulla blockchain non compare alcun dato sanitario né alcun dato che identifichi direttamente una persona: solo indirizzi pseudonimi, finalità e scadenze dei consensi, esiti degli accessi e impronte crittografiche. Referti, dettagli degli accessi e credenziali restano fuori, cifrati e sotto il controllo della struttura.

Ciò che sta su un registro distribuito è replicato su tutti i nodi e non si cancella; per questo le linee guida dell'EDPB\cite{edpb-blockchain} chiedono di non scrivere dati personali sulla catena e di permettere la cancellazione fuori da essa. Le impronte sono calcolate su dettagli che contengono un valore casuale (\emph{nonce}), quindi non si ricostruiscono provando tutte le combinazioni di medico, paziente e finalità. Il paziente compare con un indirizzo dedicato ai soli consensi, il cui legame con la persona sta nel database e può essere cancellato (sezione~\ref{sec:cancellazione}). Restano visibili alcuni metadati, discussi tra i rischi residui (sezione~\ref{sec:residui}).

\subsection{Flusso di accesso a una cartella}\label{sec:flusso}
Un accesso del medico alla cartella di un paziente attraversa quattro controlli, nell'ordine:
\begin{enumerate}
  \item \textbf{autenticazione}: il medico accede con la propria passkey e ottiene una sessione (sezione~\ref{sec:fido2});
  \item \textbf{ruolo e consenso}: il backend verifica che chi chiede abbia il ruolo di medico, assegnato dalla struttura, e che il paziente gli abbia concesso, sul registro on-chain, un consenso ancora valido per la finalità dichiarata;
  \item \textbf{audit}: l'esito, concesso o negato, viene scritto sul registro di audit; se la scrittura non riesce, nulla viene consegnato;
  \item \textbf{consegna}: solo a questo punto i referti vengono decifrati in memoria e consegnati.
\end{enumerate}
Per le situazioni d'emergenza esiste un percorso separato, che salta il controllo del consenso ma lascia una traccia distinta nell'audit (sezione~\ref{sec:accesso}). Una prova di abilitazione a conoscenza zero, descritta come estensione nella sezione~\ref{sec:zkp}, si inserirebbe tra il primo e il secondo controllo.

\subsection{Ambiente e gestione delle dipendenze}
La versione di Node.js è fissata in \texttt{.nvmrc}, quelle delle librerie nei \texttt{package-lock.json}, quella del database in \texttt{docker-compose.yml}: i risultati sono riproducibili e un aggiornamento automatico non può introdurre una dipendenza compromessa. L'immagine del database è indicata con la versione esatta, la 8.0.38, e con la sua impronta (\emph{digest}), così Docker usa sempre esattamente quella immagine. La serie~8.0 è supportata dal produttore fino a ottobre 2029 e riceve correzioni di sicurezza\cite{mongodb-lifecycle}, mentre le serie intermedie hanno un supporto molto più breve. Le versioni~8.x non si avviano sui kernel Linux dalla 6.19 alla 7.0.13, per un'incompatibilità nota della libreria di allocazione della memoria\cite{mongodb-server-121912}: serve quindi un kernel precedente al 6.19 o almeno il 7.0.14.

In esecuzione servono MongoDB in Docker, un nodo Hardhat che simula la blockchain, il backend e il frontend; i contratti sono pubblicati con un modulo di Hardhat Ignition (sezione~\ref{sec:collegamento}).

I segreti non compaiono nel codice né nel repository: credenziali del database, chiave del backend per la blockchain e indirizzi dei contratti stanno nel file \texttt{.env}, escluso da Git (il repository contiene solo il modello \texttt{.env.example}); le chiavi dell'archivio stanno nella cartella \texttt{secrets/}, anch'essa esclusa, cifrate con una passphrase che non è salvata da nessuna parte (sezione~\ref{sec:chiavi}).
